(1) 薄膜中的往返传播路程为 \(2t\)，相位差为 \(\Delta\phi_{\rm prop}=4\pi n_ft/\lambda_0\)。从低折射率向高折射率反射时电场反射系数为负，会额外产生 \(\pi\) 相位；从高向低反射则反射系数为正，不增加此相位翻转，故总干涉还须加两界面反射相位之差。
(2) 当 \(n_0<n_f<n_s\)，入射介质至薄膜及薄膜至基底两处反射都由低折射率向高折射率，均翻转 \(\pi\)，相对反射相位贡献相消。故两束反射相消要求 \(4\pi n_ft/\lambda_0=(2m+1)\pi\)，即 \(2n_ft=(m+1/2)\lambda_0\)。
(3) 法向入射的电场反射振幅为 \(r_{01}=(n_0-n_f)/(n_0+n_f)\)、\(r_{12}=(n_f-n_s)/(n_f+n_s)\)。对于理想无耗膜，在四分之一波相位厚度下严格反射振幅 \(r=(r_{01}-r_{12})/(1-r_{01}r_{12})\)；零反射要求 \(r_{01}=r_{12}\)，代数整理为 \(n_f^2=n_0n_s\)。
(4) 选设计真空波长 \(\lambda_0\)，最薄正膜厚 \(t=\lambda_0/(4n_f)\)，且 \(n_f=\sqrt{n_0n_s}\) 实现理论零反射。反射系数较大、膜层多层或宽波段时不能只把两束简单相加；精确结果应使用多重反射的几何级数或光学传输矩阵。
